\documentclass[12pt,a4paper]{article}
\usepackage[a4paper,margin=1.5cm]{geometry}
\usepackage{xcolor,listings}

\title{CSE 237C Project 1 FIR}
\author{Mitchell Bizzigotti, James Han, Benjamin Scott, Abhay Lal}

\begin{document}
\maketitle


\section*{Question 1: FIR11 Baseline}
% Implement a functionally correct 11-tap FIR filter. Do not apply pragmas or other optimizations. You can take reference from the textbook.

\begin{enumerate}
    \item[(a)] \textbf{Report latency, initiation interval (II), in clock cycles.}
    
        Latency = 19 cycles, Interval=20 cycles
    
    \item[(b)] \textbf{Report the number of BRAM, DSP, LUT and FFs used.}

        0 BRAM used, 2 DSP units used, 383 LUTs used, 733 FFs used.
\end{enumerate}


\section*{Question 2: Variable Bitwidths}
% You can specify a precise data type for each variable in your design.
% There many different data types including floating point, integer, fixed point, all with varying bitwidths and options.
% The data type provides a tradeoff between accuracy, resource usage, and performance.

Change the bitwidth of the variables inside the function body (namely \verb|coef_t| and \verb|acc_t|) using the \verb|ap_int<>| datatype.

\begin{enumerate}
    \item[(a)] \textbf{Try 3 combinations of variable bitwidths of your choice. Report the latency, II, and resource usage (BRAM, DSP, LUT, and FF) for each design in a table.}
    
        [TODO: format as table]
    
    \item[(b)] \textbf{What is the minimum bitwidth of \texttt{coef\_t} and \texttt{acc\_t} you can use without losing accuracy (i.e., your results still match the golden output)?}
    
        16 bits

\end{enumerate}


\section*{Question 3: Pipelining}
Pipelining increases throughput, typically at the cost of additional resources.
The initiation interval (II) is an important design parameter that affects performance and resource usage.

\begin{enumerate}
    \item[(a)] \textbf{Report the latency, II, and resource usage of your baseline FIR128 design.}

        Latency = 133 cycles, II = 128 cycles

        BRAM = 3, DSP = 2, FF = 610, LUT = 305
    
    \item[(b)] \textbf{Turn off the automatic pipelining by using} \verb|#pragma HLS pipeline off| \textbf{. Report the latency, II, and resource usage of this design.}

        Latency = 133 cycles, II = 128 cycles 

        BRAM = 3, DSP = 2, FF = 610, LUT = 305

    \item[(c)] \textbf{Manually pipeline the design using} \verb|#pragma HLS pipeline II=<value>| \textbf{. Report the latency, II, and resource usage of each design in a table.}

        [TODO: format as table]

    \item[(d)] \textbf{At some point setting the II to a larger value does not make sense. What is that value in this example?}

    \item[(e)] \textbf{Vitis HLS may automatically pipeline a loop without any specific pragma. Based on your observation, what is the default II for a pipelined loop used in this case?}

\end{enumerate}


\section*{Question 4: Removing Conditional Statements}
% If/else statements and other conditionals can limit the possible parallelism and often require additional resources.
% Rewriting the code to remove them can make the resulting design more efficient.

\begin{enumerate}
    \item[(a)] \textbf{Report the numbers in a table: Compare the latency, II, and resource usage of the automatically pipelined design with / without conditional statements.}
    
    \item[(b)] \textbf{Report the numbers in a table: Compare the latency, II, and resource usage of the non-pipelined design with / without conditional statements.}
    
\end{enumerate}


\section*{Question 5: Loop Partitioning}

Dividing the loop into two or more separate loops may allow for each of those loops to be executed in parallel (via unrolling), enable loop-level pipelining, or provide other benefits.
% This may increase performance and resource usage.

\begin{enumerate}
    \item[(a)] \textbf{Briefly describe the opportunity for loop partitioning in FIR128. Re-write the code to apply your idea.}
    \item[(b)] \textbf{Compare the latency, II, and resource usage of the design with / without loop partitioning.}
    \item[(c)] \textbf{Apply loop unrolling to the design with loop partitioning. Report the latency, II, and resource usage of this design.}
    \item[(d)] \textbf{What is the relationship between loop unrolling and pipelining? Can they be applied together and benefit the design? Justify you answer with experiments / references to past questions.}
\end{enumerate}


\section*{Question 6: Memory Partitioning}

The storage of the arrays in memory plays an important role in area and performance.
On one hand, you could put an array entirely in one memory (e.g., BRAM).
But this limits the number of read and write accesses per cycle.
Or you can divide the array into two or more memories to increase the number of ports.
Or you could instantiate each variable as a register allowing simultaneous access to all the variables at every clock cycle.

 Read the textbook about the memory partitioning parameters: block, cyclic, and complete.

\begin{enumerate}
    \item[(a)] \textbf{Explore array partitioning options for both arrays in your design from Question 5. Report the latency, II and resource usage. Which partition gives the best performance?}
    \item[(b)] \textbf{Loop unrolling and memory partitioning are often used together. Try disabling loop unrolling or array partitioning. Report the effects.}
\end{enumerate}


\section*{Question 7: Best Design}

Combine any number of optimizations to get your best architecture.
A design with high throughput will likely take a lot of resources.
A design with small resource usage likely will have lower performance, but that could still be good enough depending the application goals.

\begin{enumerate}
    \item[(b)] \textbf{Report the resource usage of your design with the best throughput. Explain why the resource usage is high compared with the baseline in Question 2.}
    \item[(a)] \textbf{Combine any number of optimizations to get your best throughput. What optimizations did you use to obtain this result? Report the latency, II, throughput (in MHz). It is possible to create a design that outputs a result every cycle, i.e., get one sample per cycle, so a throughput of 100 MHz (assuming a 10 ns clock).}
\end{enumerate}

\end{document}
